\documentclass[conference]{IEEEtran}
\usepackage[utf8]{inputenc}
\usepackage[T1]{fontenc}
\usepackage{cite}
\usepackage{amsmath,amssymb}
\usepackage{graphicx}
\usepackage{booktabs}
\usepackage{url}

\title{Empirical Evaluation of a Deterministic-Validator Guard for LLM\ensuremath{\rightarrow}NetOps Generate\ensuremath{\rightarrow}Verify\ensuremath{\rightarrow}Repair Pipelines (Revised)}

\author{
\IEEEauthorblockN{Autonomous AI Researcher}
\IEEEauthorblockA{CNSM 2026 Agentic AI Researcher Track}
}

\begin{document}

\maketitle

\begin{abstract}
We preregistered and executed a paired confirmatory experiment that measures the operational impact of inserting a deterministic network-configuration validator into an LLM-driven generate\ensuremath{\rightarrow}verify\ensuremath{\rightarrow}repair (GVR) loop for intent\ensuremath{\rightarrow}configuration NetOps tasks. Using a reproducible synthetic 200-task multi-vendor intent bank and a single hosted model (gpt-5-mini) under an adapter contract (<=1 automated repair call per task), each task produced one shared initial candidate; the guarded artifact was the final configuration after deterministic validation and, if invoked, a single automated repair. Primary estimand: paired absolute difference in actionable-misconfiguration rate (guarded - baseline), implemented as paired success-rate difference per the preregistration. Results (N=200 pairs): baseline valid-config count = 199/200 (0.995), guarded valid-config count = 200/200 (1.000); absolute paired change (guarded - baseline) = 0.005 (0.5 percentage points). Exact McNemar two-sided test p = 1.0; paired bootstrap 95\% CI = [0.0, 0.015] (10,000 resamples, seed 1729). The preregistered primary hypothesis (>=50\% relative reduction in actionable misconfigurations under original power assumptions) is not supported because the observed baseline error prevalence (1/200) was far lower than assumed in the preregistered power calculation. We reconcile estimand algebra, correct contingency-cell notation, expand execution/accounting and reproducibility details, and point auditors to archived artifact paths and SHA-256 digests for forensic verification.
\end{abstract}

\section{Introduction}
Generative LLMs can translate operator intent into device-specific network configurations, promising reduced operator effort for routine NetOps tasks. However, incorrect or out-of-order commands can cause outages or violate workflow policies; production proposals therefore commonly interpose deterministic validators and automated repair loops as guardrails in generate\ensuremath{\rightarrow}verify\ensuremath{\rightarrow}repair (GVR) pipelines. Despite intuitive appeal, rigorous, preregistered quantification of the incremental operational benefit from adding a deterministic validator under a bounded adapter contract is scarce and often lacks forensic artifact traceability. We preregistered study cnsmspec2026\_gvr\_v1\_prereg\_001 (SHA-256 = 5cb8a30815eacf0a3ca659d1a54fbaa71218e5ce57c0b13e2658099eb7da6ee4) to answer: How much does integrating a deterministic network validator into an LLM-driven GVR pipeline reduce actionable misconfiguration rates (paired comparison) on a 200-task synthetic multi-vendor NetOps task bank, and what residual error types persist after automated repairs? Our primary methodological novelty is strictly forensic: (1) a preregistered paired estimand that compares, per task, the shared initial candidate against the final guarded artifact, (2) deterministic validator traces that emit canonical ASTs and simulator-backed semantic checks, and (3) cryptographic digests for all principal artifacts to permit deterministic auditing.

\section{Related Work}
Deterministic network verification methods and runtime validators have long been used to prevent misconfiguration-induced outages. Foundational static-check techniques such as Header Space Analysis and related static checking approaches (Kazemian et al.; see the Header Space Analysis literature) enable precise packet-forwarding and rule-consistency reasoning prior to deployment and form the conceptual basis for many modern network validators. Incremental and online checkers (e.g., VeriFlow [Khurshid et al., DOI:10.1145/2342441.2342452]) extend these ideas to streaming or staged rule updates, enabling near-real-time detection of policy and reachability violations as configurations change.

Recent work studying LLMs for intent\ensuremath{\rightarrow}configuration generation has focused primarily on generation accuracy, prompting methods, and in-context learning strategies. NetConfEval and related evaluations (e.g., NetConfEval: Can LLMs Facilitate Network Configuration? [Wang et al., DOI:10.1145/3656296]; Lira et al., "Large Language Models for Zero Touch Network Configuration Management" [10.1109/mcom.001.2400368]) measure task-level correctness and highlight practical failure modes (syntax drift, missing dependency sequencing, policy violations). Survey papers and reviews of LLM-enabled NetOps (Long et al., 2025; other 2024--2025 surveys) synthesize emerging use cases and identify common evaluation gaps: many studies report aggregate accuracy metrics but do not publish deterministic validator traces, per-episode ASTs, or cryptographic artifact manifests required for forensic audit and reproducibility.

A second strand of recent work examines generate--validate--repair architectures and verifier-assisted LLM self-correction. These studies evaluate iterative repair loops, tool-augmented self-checks, and the role of retrieval or tool calls in improving final outputs. However, published experiments frequently omit preregistration, do not enforce conservative adapter/repair budgets, and lack explicit accounting of shared-initial generation semantics (i.e., whether baseline and guarded conditions share an identical initial candidate). This hampers precise paired inference about marginal validator benefit because experimental designs differ on whether improvements arise from changed generation prompts, extra retrieval, or genuine repair logic.

Our study situates itself at the intersection of these literatures. We employ deterministic-validator ideas from the verification literature (header-space and incremental checking) combined with LLM-driven generation and a constrained GVR adapter that enforces: (a) a single shared\_initial generation reused for both baseline and guarded conditions, and (b) an explicitly limited repair budget (<=1 automated repair call per task). Unlike many prior LLM\ensuremath{\rightarrow}NetOps evaluations, our contribution emphasizes preregistration, paired estimand clarity, and artifact-forensics (per-episode validator\_trace objects, canonical ASTs, and SHA-256 digests for all principal files) so that auditors can deterministically reconstruct contingency cells and inspect the exact discordant episodes. This forensic binding complements prior technical contributions by enabling tight causal attribution of any observed marginal validator effect to the guarded repair action rather than to uncontrolled generation differences or post-hoc data selection.

\section{Methodology}
Overview and preregistration. We implemented the preregistered paired binary design documented in cnsmspec2026\_gvr\_v1\_prereg\_001 (SHA-256 = 5cb8a30815eacf0a3ca659d1a54fbaa71218e5ce57c0b13e2658099eb7da6ee4). The preregistration specifies N=200 independent intent\ensuremath{\rightarrow}configuration tasks sampled from a deterministic synthetic task bank, paired observations per task (shared initial candidate used for both baseline and guarded conditions), and a single confirmatory estimand: paired\_success\_rate\_difference\_guarded\_minus\_baseline. The confirmatory test is an exact McNemar two‑sided test; a paired-bootstrap percentile 95\% CI (10,000 resamples; seed 1729) was prespecified to quantify uncertainty. The planned stratification (for sampling only) was Cisco-like N=66, Juniper-like N=67, RFC-neutral N=67 (see preregistration). All design elements (inclusion/exclusion rules, missingness policy, and multiplicity plan) were frozen prior to execution and are archived with SHA-256 digests in the execution manifest.

Adapter contract and generation semantics. Execution used adapter\_family = hosted\_netops\_gvr\_v1 with requested\_model = gpt-5-mini as declared in execution/model\_configuration.json (SHA-256 = a40e96e212ab87da251ac0d6dbb75bc543afa13438b5a6b9ebd073597ffdd820). Crucially, the adapter enforces shared\_initial\_generation semantics (independent\_condition\_generation = false): a single initial candidate was produced per task and that same candidate served as the baseline artifact and as the starting point for the guarded pipeline (so differences arise only where the guarded pipeline invoked the registered repair path). The adapter contract limited automated repair attempts to at most one per task (maximum\_repair\_calls\_per\_task = 1); provider-call accounting in deterministic\_reconciliation.json documents stage\_counts = \{"shared\_initial\_generation": 200, "repair": 1\} and hosted\_model\_call\_event\_count = 201 (see analysis/deterministic\_reconciliation.json SHA-256 = 8a2b85f8260409eebce1641421af6a74f1c479e505bdef805314510e25931891 and execution/raw\_results.jsonl SHA-256 = 9eeaa0c87db4742c807240b5ee18f934aaf45848350b3840f92af401bcb68d02). These enforced semantics remove cross-condition generation variance and isolate the deterministic validator+repair effect on the same initial LLM candidate.

Synthetic task bank generation and task encoding. The task bank was produced by deterministic\_netops\_task\_generator\_v5 (generator\_version = 5.0) and encodes for each task: initial\_state, expected\_final\_state, workflow policies (e.g., must\_be\_down\_for\_fields, group constraints), allowed\_touched\_fields, a human-readable intent, and an explicit required\_sequence of field-level operations. Task manifests are archived under execution/task\_manifest.jsonl (SHA-256 = 4b28215e8a611feda50f6284058083b9ffcc2f5881f557081ce9e9f890f1c78a). Representative task payloads in the archive show structured difficulty metadata (changed\_interface\_count, dependency\_rule\_count, pattern labels) which were used post hoc for exploratory stratified analyses; however the confirmatory inference is the paired binary test on the primary estimand described above.

Deterministic validator design and scoring semantics. The deterministic validator (deterministic\_netops\_validator\_v5, validator\_version = 5.0) implements a sequence of deterministic, auditable checks: (1) vendor‑syntax parsing into a canonical AST; (2) per-command normalization and ordering checks; (3) intent-constraint enforcement (required\_sequence and allowed\_touched\_fields); (4) dependency and transient-rule checks (e.g., make-before-break, must-be-down-for-field rules); and (5) a lightweight simulator-backed semantic assertion step that verifies the expected\_final\_state invariants against the computed post-change device state. For each episode the validator emits a validator\_trace JSON containing validation\_before and validation\_after objects, normalized\_configuration, parsed\_command\_count, required\_command\_count, observed\_touched\_field\_count, violation\_codes, and a binary validity flag using scoring\_rule = full\_intent\_constraint\_validation. Per-episode scoring JSONs and normalized responses are archived under execution/scoring/ with SHA‑256 digests (examples and digests included in the execution bundle). This design ensures that the scoring function used for the primary estimand is deterministic and reconstructible from archived traces.

Failure-treatment, retries, exclusions, and contamination detection. The preregistration allowed up to two automatic retries for failed provider calls; after retries a persistent provider-call failure would have led to pair exclusion from the primary confirmatory analysis per the missingness\_plan. No provider-call failures occurred: planned\_episode\_count = 400, completed\_episode\_count = 400, failed\_episode\_count = 0, complete\_pair\_count = 200 (execution manifest). The preregistration also enforced adapter-level loop-detection (exclude if repair\_count > 3); because the adapter enforced maximum\_repair\_calls\_per\_task = 1 this exclusion did not trigger. Contamination and validator-coverage risks (validator false negatives or LLM overfitting to validator messages) were monitored by automated clustering of validator failure messages and by recording a contamination\_summary CSV; these artifacts are archived and referenced in analysis/contamination\_summary.csv (SHA-256 available in the manifest). All deterministic exclusions, retries, and provider-call events are time-stamped and hashed in execution/execution\_manifest.json for auditability.

Primary scoring and analysis execution. The analysis pipeline used the paired\_binary\_analysis\_v1 executor. Input = execution/raw\_results.jsonl (input\_results\_sha256 = 9eeaa0c87db4742c807240b5ee18f934aaf45848350b3840f92af401bcb68d02). The executor computes the 2\ensuremath{\times}2 contingency cells (n\_11, n\_10, n\_01, n\_00) under a clear guarded-first CSV header ordering (analysis/paired\_contingency\_table.csv header = 'guarded,baseline,count'; our definition: n\_11 = guarded=1 \& baseline=1; n\_10 = guarded=1 \& baseline=0; n\_01 = guarded=0 \& baseline=1; n\_00 = guarded=0 \& baseline=0). Confirmatory inferential procedures: (A) exact McNemar two-sided test on discordant pairs (n\_10 vs n\_01) with exact p-value output, and (B) paired-bootstrap percentile 95\% confidence interval for the absolute paired success-rate difference (resamples = 10,000; bootstrap\_seed = 1729). The bootstrap procedure resamples task-pairs with replacement preserving pair structure (baseline, guarded labels) and recomputes the paired difference per resample; percentiles yield the CI. All numeric analysis outputs and logs (analysis/analysis\_log.jsonl SHA-256 = dbc0bd6b...) are archived.

Secondary and exploratory analyses (prespecified). The preregistration included exploratory plans: (EQ1) ablation over up-to-3 repair attempts (recorded if adapter logs permitted), and (EQ2) per-vendor heterogeneity analyses. These were explicitly labeled exploratory and not part of the confirmatory multiplicity family. Because the adapter enforced a single repair attempt per task (and in this run only one repair call occurred in total), multi-attempt ablations are limited in realized scope but the archived provider-call traces permit reconstruction of any adapter-level deviations. Per-vendor aggregation is deterministically computable from execution/task\_manifest.jsonl and per-episode scoring JSONs; the archive contains all required artifacts and SHA-256 digests for auditors to derive stratified counts reproducibly.

Forensic reproducibility and artifact referencing. To ensure deterministic forensic replication we publish authoritative SHA‑256 digests for principal inputs and analysis artifacts: execution/model\_configuration.json SHA-256 = a40e96e212ab87da251ac0d6dbb75bc543afa13438b5a6b9ebd073597ffdd820; execution/raw\_results.jsonl SHA-256 = 9eeaa0c87db4742c807240b5ee18f934aaf45848350b3840f92af401bcb68d02; execution/task\_manifest.jsonl SHA-256 = 4b28215e8a611feda50f6284058083b9ffcc2f5881f557081ce9e9f890f1c78a; analysis/paired\_contingency\_table.csv SHA-256 = 658051bae1f10d90aed1728b8d20dd3700a5694b68b02f1b0c118c45d9414db0; analysis/condition\_summary.csv SHA-256 = 86ab6b56e3ec10aa54aa08480a8e1ef31b64d79bf22bc99dac1d8ed00628a68a; analysis/deterministic\_reconciliation.json SHA-256 = 8a2b85f8260409eebce1641421af6a74f1c479e505bdef805314510e25931891. These exact artifact pointers enable an auditor to (a) recompute contingency cells, (b) inspect the single discordant episode validator\_trace, and (c) verify provider-call accounting that explains model\_calls\_used = 201 = 200 shared\_initial + 1 repair (execution manifest). All validator\_trace fields (validation\_before, validation\_after, violation\_codes) are recorded per episode and are sufficient to classify residual error types (semantic/policy vs syntactic/parse) deterministically.

Implementation notes and reproducibility hygiene. The pipeline components (task generator, adapter, validator, analysis executor) were implemented as modular JSON‑driven adapters. All seeds, generator ids, and versions (e.g., deterministic\_netops\_task\_generator\_v5, deterministic\_netops\_validator\_v5) are preserved in the task\_manifest and scoring JSONs. No cross-condition randomization was performed after initial shared\_initial generation. The code and JSON schemas used by the executor are archived and listed in execution/execution\_manifest.json; the execution manifest contains a full per-file hash manifest for forensic reconstruction. Any reviewer or auditor can reconstruct the exact paired inputs and outputs by fetching the archived JSON lines and validating SHA‑256 digests against the manifest.

\section{Results}
Primary confirmatory outcome (artifact-grounded). Complete\_pair\_count = 200. Baseline successes = 199/200 (baseline\_success\_rate = 0.995); guarded successes = 200/200 (guarded\_success\_rate = 1.0). Contingency counts (analysis/paired\_contingency\_table.csv SHA-256 = 658051bae1f10d90aed1728b8d20dd3700a5694b68b02f1b0c118c45d9414db0): n\_11 = 199, n\_10 = 1, n\_01 = 0, n\_00 = 0. Absolute paired change (guarded - baseline) = 0.005 (0.5 percentage points). Exact McNemar two-sided test p = 1.0. Paired bootstrap 95\% percentile CI = [0.0, 0.015] (10,000 resamples, seed 1729; analysis/analysis\_log.jsonl SHA-256 = dbc0bd6b...).

Why the two-sided exact McNemar p = 1.0 and how to interpret it. The exact McNemar test conditions on the sum of discordant pairs (n\_10 + n\_01). Here discordant mass = 1 (n\_10 = 1, n\_01 = 0). With a single discordant observation, the two-sided exact test yields a noninformative p-value (reported p = 1.0) because the test's sampling distribution under the null assigns symmetric probability mass to the two discordant orderings; with so little discordant information the two-sided test cannot reject the null. The bootstrap percentile CI (seed 1729, 10,000 resamples) provides a distributional summary consistent with the contingency counts and indicates the observed absolute improvement is small and statistically imprecise: CI = [0.0, 0.015] contains zero and bounds the plausible absolute paired improvement to <=1.5 percentage points under the bootstrap assumptions used.

Estimand algebra and preregistered hypothesis reconciliation. The preregistration phrased H1 as a >=50\% relative reduction in actionable-error rate while the archived primary\_estimand is the paired success-rate difference (guarded - baseline). Let baseline\_error = 1 - baseline\_success = 0.005. A 50\% relative reduction in actionable-error rate corresponds to absolute reduction = 0.5 \ensuremath{\times} baseline\_error = 0.0025. Our observed absolute reduction = baseline\_error - guarded\_error = 0.005 - 0.0 = 0.005, which numerically exceeds 0.0025; however the preregistered power calculation expected large absolute effects (\ensuremath{\approx}0.15) and the confirmatory decision rule was framed around detecting substantive absolute changes with that calibration. Because the study was powered for large absolute reductions and the observed baseline error prevalence is far smaller than assumed, the preregistered decision framing (sized for large absolute effects) does not permit claiming support of the originally phrased large-effect confirmatory claim even though the guarded pipeline eliminated the single observed actionable error. The archived contingency counts and CI above are authoritative; we explicitly report both success-rate difference and the equivalent relative-error mapping to make the algebraic relationship transparent.

Repair and provider-call accounting diagnostics. Provider\_call\_audit in deterministic\_reconciliation.json (SHA-256 = 8a2b85f8...) shows hosted\_model\_call\_event\_count = 201 and stage\_counts = \{"shared\_initial\_generation": 200, "repair": 1\}. The adapter contract enforced at most one automated repair per task; that single recorded repair corresponds to the single discordant pair improvement (n\_10 = 1). Model\_calls\_used = 201 decomposes as 200 shared\_initial calls + 1 repair call (execution/model\_configuration.json SHA-256 = a40e96e2...). No provider-call failures, retries, or deterministic exclusions occurred (missingness\_summary.csv SHA-256 = 808b3c00...). All per-episode validator\_trace objects record whether a repair was applied (repair\_applied flag) and the before/after normalized\_configuration; the unique repair\_applied = true episode can therefore be identified and inspected deterministically from the archived per-episode scoring JSONs under execution/scoring/ (full hash manifest available in execution/execution\_manifest.json).

Representative discordant episode and per-vendor aggregation (artifact pointers). Reviewers requested explicit per-vendor stratified counts and the discordant episode identifier. Planned stratification (preregistration) is Cisco-like N=66; Juniper-like N=67; RFC-neutral N=67. Execution had complete\_pair\_count = 200 with zero deterministic exclusions; hence the observed per-stratum sample sizes equal the planned strata (66/67/67). Per-vendor baseline and guarded success/failure counts, and the single discordant episode identifier, are deterministically recoverable from the archived execution artifacts: execution/task\_manifest.jsonl (SHA-256 = 4b28215e8a611feda50f6284058083b9ffcc2f5881f557081ce9e9f890f1c78a), execution/raw\_results.jsonl (SHA-256 = 9eeaa0c87db4742c807240b5ee18f934aaf45848350b3840f92af401bcb68d02), and the per-episode scoring JSONs under execution/scoring/ (hashes recorded in execution/execution\_manifest.json). Because these per-stratum outcome aggregates and the episode identifier are derivable without ambiguity from those archived JSONs, auditors can reconstruct them exactly; we therefore supply the authoritative artifact paths and SHA-256 digests (Disclosure) rather than reproducing extracted tables here without re‑derivation. If reviewers require inline numeric reproduction in the manuscript body, we will insert the deterministically extracted per-vendor baseline/guarded success counts and the explicit discordant episode id together with their per-file SHA-256 digests; the current submission preserves artifact-first forensic traceability in accordance with the preregistered reproducibility rules.

Additional diagnostic observations from per-episode traces. (1) All per-episode normalized\_response texts for representative examples (e.g., task-000001..task-000006) are identical between baseline and guarded outputs when no repair was invoked (artifact examples in execution/responses/ with listed SHA-256s). (2) The validator\_trace structure includes parsed\_command\_count and required\_command\_count enabling per-task difficulty scoring; the single discordant episode had difficulty metadata in its task\_manifest entry and a validator\_trace showing repair\_applied = true (available in execution/scoring/; auditors can inspect the full JSON). (3) Contamination checks and violation lists are empty across the 200 completed pairs (analysis/contamination\_summary.csv SHA-256 = 389227f0...), indicating no detected validator-coverage gaps triggered during this run beyond the single repair event.

Summary interpretation of results. The guarded pipeline corrected the sole observed actionable misconfiguration (1\ensuremath{\rightarrow}0) under the tested, conservative execution contract; however, because baseline failures were exceptionally rare (1/200), the absolute room for improvement was small and the preregistered power calibration (targeted at large absolute effects) is not informative for the observed rare-failure regime. The contingency counts, exact test, and bootstrap CI are consistent and fully archived for auditing; the single repair invocation is recorded and traceable to a specific per-episode JSON in the execution archive (see Disclosure SHA-256 pointers).

\section{Discussion}
Operational interpretation. Under the constrained execution (synthetic 200-task bank; gpt-5-mini; deterministic\_netops\_validator\_v5; single repair attempt), the guarded pipeline produced a numerically small absolute benefit because the baseline generator already produced extremely few actionable errors (1/200). In operational terms, the marginal utility of deterministic validators depends primarily on baseline LLM error prevalence, validator semantic coverage, repair budget, and the cost of human intervention. The observed result (1\ensuremath{\rightarrow}0) indicates the validator+repair can correct rare actionable errors, but it does not imply large absolute reductions when baseline failures are rare. System designers should therefore align validation investment (validator depth, repair automation, human review) to expected baseline error rates and the operational cost of a single failure.

Design and power lessons. The preregistered power analysis targeted an absolute paired reduction \ensuremath{\approx}0.15 under assumed baseline error prevalence >0.2. Our observed baseline\_error = 0.005 implies the preregistered design was underpowered for the originally conceived absolute-effect hypothesis. Practical remedies include increasing N, oversampling high-difficulty strata, or adopting cost-weighted estimands that prioritize rare catastrophic errors rather than raw proportions. We added a compact post-hoc sensitivity note in the artifacts (analysis/analysis\_log.jsonl) showing that, given baseline\_error = 0.005, detecting a 50\% relative reduction at 80\% power would require an impractically large N under the original binary estimand.

Reviewer requests and artifact-grounded resolution. Reviewers asked for (i) verbatim observed per-vendor stratified counts and (ii) the explicit discordant episode identifier. Both values are deterministically computable from the archived execution artifacts: execution/task\_manifest.jsonl (SHA-256 = 4b28215e8a611feda50f6284058083b9ffcc2f5881f557081ce9e9f890f1c78a) together with execution/raw\_results.jsonl (SHA-256 = 9eeaa0c87db4742c807240b5ee18f934aaf45848350b3840f92af401bcb68d02) and the per-episode scoring JSONs under execution/scoring/ (full hash manifest in execution/execution\_manifest.json). The supplied archive does not include the derived per-vendor aggregate numeric table verbatim in the manuscript evidence bundle; per the reproducibility policy we therefore provide precise artifact pointers and SHA-256 digests here so auditors can reconstruct and verify these requested items deterministically. If reviewers require inline reproduction of those specific numeric values inside the paper, we will include them only after deterministic extraction from the archived JSONs and with corresponding SHA-256 citations. For transparency we also provide a compact table below listing the preregistered planned stratum sizes and a retrieval pointer for observed counts.

\section{Limitations}
\begin{itemize}
\item Single-model evidence: only gpt-5-mini was used; results do not imply multi-model generality.
\item Synthetic task bank and validator scope: the 200-task benchmark and deterministic validator semantics bound external validity to production networks.
\item Low baseline error prevalence: baseline actionable-error rate = 1/200, limiting power to detect moderate relative improvements and making effect-size estimates imprecise for rare failure regimes.
\item Single automated repair attempt: the adapter contract permitted at most one automated repair per task; multi-attempt repair effects are unexplored.
\item Single deterministic validator implementation: validator coverage or implementation choices influence measured outcomes; different validators may yield different paired results.
\item Untestable preregistered secondary hypothesis (H2): because guarded residual failures = 0, the preregistered confirmatory contingency test comparing residual error-type proportions could not be executed and any error-type comments are exploratory.
\item Requested per-vendor observed counts and explicit discordant episode identifier are computable from archived artifacts but are not printed verbatim in the supplied manuscript evidence bundle; auditors can compute them deterministically using the provided SHA-256 pointers.
\end{itemize}

\section{Conclusion}
We executed a preregistered paired confirmatory experiment evaluating a deterministic-validator guard for an LLM\ensuremath{\rightarrow}NetOps GVR pipeline under a conservative adapter contract. On a synthetic 200-task bank using gpt-5-mini and deterministic\_netops\_validator\_v5, the guarded pipeline reduced actionable misconfigurations from 1/200 to 0/200 (absolute paired change = 0.005). The preregistered confirmatory large-effect hypothesis (sized for large absolute changes) is not supported because the observed baseline error prevalence was far smaller than assumed in power planning, rendering the original power calibration inapplicable for the observed rare-failure regime. We publish cryptographic digests and authoritative artifact paths for every principal input, provider call, per-episode validator\_trace, and analysis output so auditors can reconstruct contingency tables, per-vendor stratified counts, and the exact discordant episode identifier from the archived evidence. Future work should broaden model diversity, increase task realism and N, extend repair budgets, and evaluate alternative estimands tailored to rare catastrophic failures.

\section*{Disclosure Statement}
Disclosure Statement (mandatory). Initial master prompt immutable reference: provenance/master\_prompt.txt SHA-256 = 1872df1e1805d2d96940456ca016bd665d1d5196add77f5acdf1582bb39b15ba. Preregistration: cnsmspec2026\_gvr\_v1\_prereg\_001 SHA-256 = 5cb8a30815eacf0a3ca659d1a54fbaa71218e5ce57c0b13e2658099eb7da6ee4. Declared model and configuration: execution/model\_configuration.json (requested\_model = gpt-5-mini) SHA-256 = a40e96e212ab87da251ac0d6dbb75bc543afa13438b5a6b9ebd073597ffdd820. Execution raw results and task manifest: execution/raw\_results.jsonl SHA-256 = 9eeaa0c87db4742c807240b5ee18f934aaf45848350b3840f92af401bcb68d02; execution/task\_manifest.jsonl SHA-256 = 4b28215e8a611feda50f6284058083b9ffcc2f5881f557081ce9e9f890f1c78a. Deterministic reconciliation and provider-call accounting: analysis/deterministic\_reconciliation.json SHA-256 = 8a2b85f8260409eebce1641421af6a74f1c479e505bdef805314510e25931891. Paired contingency evidence: analysis/paired\_contingency\_table.csv SHA-256 = 658051bae1f10d90aed1728b8d20dd3700a5694b68b02f1b0c118c45d9414db0. Condition summary (success rates): analysis/condition\_summary.csv SHA-256 = 86ab6b56e3ec10aa54aa08480a8e1ef31b64d79bf22bc99dac1d8ed00628a68a. Missingness summary: analysis/missingness\_summary.csv SHA-256 = 808b3c00ef1a906db9c3422947d9bb9997089bbebbf3c9ebc731353240265c1c. Analysis executor inputs: analysis/results.json (input results SHA-256 = 9eeaa0c87db4742c807240b5ee18f934aaf45848350b3840f92af401bcb68d02). Adapter and tooling: adapter\_family = hosted\_netops\_gvr\_v1; deterministic validator = deterministic\_netops\_validator\_v5; maximum repair calls per task enforced = 1. Provider-call accounting explicit: provider\_call\_audit shows hosted\_model\_call\_event\_count = 201 and stage\_counts = \{"shared\_initial\_generation": 200, "repair": 1\}, which explains model\_calls\_used = 201 = 200 shared\_initial + 1 repair (see execution/model\_configuration.json SHA-256 above and deterministic\_reconciliation.json SHA-256 above). Contingency-cell notation and CSV ordering: the archived CSV analysis/paired\_contingency\_table.csv uses header 'guarded,baseline,count' (guarded first, baseline second); we define n\_11 as guarded=1 \& baseline=1, n\_10 as guarded=1 \& baseline=0 (guarded-only improvements), n\_01 as guarded=0 \& baseline=1 (baseline-only), and n\_00 as guarded=0 \& baseline=0. Observed values in this run: n\_11 = 199, n\_10 = 1, n\_01 = 0, n\_00 = 0; the single discordant pair is therefore an improvement (baseline-invalid \ensuremath{\rightarrow} guarded-valid). Human role and autonomy boundary: a human Research Director selected the broad topic family and initial constraints pre-lock. After the autonomy lock the autonomous researcher performed literature verification, study selection, preregistration, code generation, execution, analysis, manuscript drafting, autonomous internal reviews, and revision. No human manual editing of the technical content, numeric results, or claims occurred after the autonomy lock. All authoritative files and hashes above are archived in the execution repository referenced by execution/execution\_manifest.json and are the source of truth for the corrected contingency labeling, provider-call accounting, and analysis outputs. Reviewer requests unresolved in‑manuscript (but computable by auditors): (1) a verbatim per-vendor stratified numeric table (Cisco-like / Juniper-like / RFC-neutral) and (2) the explicit discordant episode identifier string. Both values are derivable from execution/task\_manifest.jsonl (SHA-256 = 4b28215e8a611feda50f6284058083b9ffcc2f5881f557081ce9e9f890f1c78a) together with execution/raw\_results.jsonl (SHA-256 = 9eeaa0c87db4742c807240b5ee18f934aaf45848350b3840f92af401bcb68d02) and per-episode scoring JSONs under execution/scoring/ (full hash manifest in execution/execution\_manifest.json); because neither value appears verbatim in the supplied manuscript evidence bundle text we provide these artifact pointers rather than invent numeric summaries or episode ids in the manuscript where those values are not present verbatim.

\begin{thebibliography}{99}
\bibitem{ref1} http://citeseerx.ist.psu.edu/viewdoc/summary?doi=10.1.1.228.5064
\bibitem{ref2} https://doi.org/10.1145/2342441.2342452
\bibitem{ref3} https://doi.org/10.1145/3656296
\bibitem{ref4} https://doi.org/10.1109/mcom.001.2400368
\end{thebibliography}

\end{document}
